\section{The arithmetic trace on a finite set of local cochains}
\label{sec:do27-trace}

An arithmetic class of degree three can be evaluated through local
classes of degree two. The passage uses a gluing complex and a choice
of primitives. Its sign is forced by the differential. We construct
this passage and compare it with the arithmetic orientation.

Retain the totally imaginary field \(F\), the ordinary small site of
\(X=\operatorname{Spec}\mathcal O_F\), the positive integer \(n\),
and the chosen primitive root \(\zeta\) from
Section~\ref{sec:do26-orientation}. Put \(M=\underline{\mathbb Z/n}\).
Let \(S\) be a nonempty finite set of finite primes containing every
prime above \(n\), and write
\[
 U=X\setminus S,\qquad
 R_v=\mathcal O_{X,v}^{h},\quad D_v=\operatorname{Spec}R_v,
 \quad K_v=\operatorname{Frac}R_v.
\]
The completion of \(K_v\) is \(F_v\), with valuation ring
\(\widehat R_v\). All sites remain ordinary small \'etale sites.
Every scheme map below gives a morphism of these sites by base change.
Cohomological degrees increase under the differential. Coefficient
sheaves and Galois modules are discrete, and products indexed by \(S\)
are finite. The complexes are algebraic complexes of abelian groups.
Sheaf tensor products are over \(\underline{\mathbb Z}\); none is
completed. We write multiplicative abelian groups additively in complexes.

Use the coefficient homomorphism
\[
 \epsilon=\iota e_\zeta:M\longrightarrow\mathbb G_m,
 \qquad a\longmapsto\zeta^a,
\]
where \(\iota:\mu_{n,X}\hookrightarrow\mathbb G_m\).
Lemma~\ref{lem:do26-roots} identifies its kernel-valued factor with
the constant sheaf, including at coefficient primes. The nonzero
power-map cokernel in Lemma~\ref{lem:coefficient-prime-defect} remains
part of the coefficient sequence on \(X\).

\subsection{The arithmetic input and the local groups}

We use the class-field trace and divisor normalization of
\eqref{eq:do26-trace-input}--\eqref{eq:do26-trace-normalization}.
For the generic point of a henselian discrete valuation ring with finite
residue field, the same local class-field input gives
\(R^{>0}j_*\mathbb G_m=0\).
Its stated scope includes henselian, as well as complete, rings
\cite[\S1, (1.1), p.~525]{Mazur1973}. In particular, with
\(\delta_S\) the divisor connecting map for \(S\),
\begin{equation}
\label{eq:do27-input}
 \operatorname{tr}_F\delta_S((t_v)_{v\in S})=\sum_{v\in S}t_v.
\end{equation}
The connecting map lifts a cocycle and applies the differential.
This sign and the arithmetic Frobenius normalization are fixed before
the gluing complex is formed.

\begin{lemma}[The groups needed at a boundary prime]
\label{lem:do27-local}
For \(v\in S\), one has
\[
 H^q(D_v,\mathbb G_m)=0\ (q>0),\qquad
 H^q(D_v,M)=0\ (q\geq2).
\]
The same statements hold for \(\operatorname{Spec}\widehat R_v\).
On either fraction field \(L=K_v,F_v\),
\[
 H^1(L,\mathbb G_m)=0,\quad
 \operatorname{inv}_v:H^2(L,\mathbb G_m)\xrightarrow{\sim}\mathbb Q/\mathbb Z,
 \quad H^q(L,\mathbb G_m)=0\ (q\geq3).
\]
Consequently \(\epsilon\) gives
\[
 H^2(L,M)\xrightarrow{\sim}\tfrac1n\mathbb Z/\mathbb Z,
 \qquad H^3(L,M)=0.
\]
Completion preserves the displayed invariants.
\end{lemma}

\begin{proof}
Let \(R\) be either valuation ring, with residue field \(k\).
A strictly henselian local scheme has exact global sections on its
small \'etale site. Indeed, any \'etale covering has a member with a
point above the closed point. Hensel's lemma supplies a section of
that member. This lifts each section through a sheaf epimorphism.
Finite \'etale neighborhoods of the closed point of
\(\operatorname{Spec}R\) are cofinal among its \'etale neighborhoods.
Descent along their strict henselian limit therefore computes
\[
 H^q(R,\mathcal F)=H^q_{\mathrm{cont}}
       (G_k,\mathcal F(R^{\mathrm{sh}})).
\]
For \(\mathbb G_m\), the module is
\(E=(R^{\mathrm{sh}})^\times\). The unit-norm and cyclic-resolution
arguments of Lemmas~\ref{lem:do26-norm} and
\ref{lem:do26-local-invariant} give \(H^{>0}(G_k,E)=0\).
They apply to either valuation ring and to unramified extensions of
arbitrary degree, including degrees divisible by the residue
characteristic. For \(M\), the cyclic-resolution and inflation
argument in Lemma~\ref{lem:do26-defect} gives
\(H^q(G_k,M)=0\) for \(q\geq2\), also at coefficient primes.

The local direct-image input identifies fraction-field cohomology
with the cohomology of \(G_k\) acting on
\(L^{\mathrm{nr},\times}\). The valuation sequence is
\[
 0\longrightarrow E\longrightarrow L^{\mathrm{nr},\times}
   \xrightarrow{v}\mathbb Z\longrightarrow0.
\]
Valuation is surjective on invariants, because the base field has a
uniformizer. Thus degree one vanishes. In degree two the sequence
identifies the group with \(H^2(G_k,\mathbb Z)\).
The sequence \(0\to\mathbb Z\to\mathbb Q\to\mathbb Q/\mathbb Z\to0\)
then identifies it with \(H^1(G_k,\mathbb Q/\mathbb Z)\).
Positive cohomology of \(\mathbb Q\) vanishes by averaging on finite
quotients. Finite torsion cohomology of \(G_k\) vanishes above degree
one, so the higher groups in the assertion vanish.
Evaluation at arithmetic Frobenius defines \(\operatorname{inv}_v\).

The fraction field has characteristic zero. Its \'etale Kummer sequence
is therefore exact. The degree-two assertion follows from the vanishing
of \(H^1(L,\mathbb G_m)\). Multiplication by \(n\) on
\(H^2(L,\mathbb G_m)=\mathbb Q/\mathbb Z\) is surjective, so the
same sequence gives \(H^3(L,M)=0\). The chosen root identifies
\(M\) with \(\mu_n\) on this field.

For an unramified cyclic extension of degree \(d\), let
\(\sigma\) be arithmetic Frobenius and \(\varpi\) a uniformizer.
The cocycle of Lemma~\ref{lem:do26-local-invariant} is
\begin{equation}
\label{eq:do27-carry}
 \alpha(\sigma^i,\sigma^j)=
 \varpi^{\lfloor(i+j)/d\rfloor},\qquad 0\leq i,j<d,
 \qquad \operatorname{inv}_v[\alpha]=1/d.
\end{equation}
Completion preserves the unramified extension, arithmetic Frobenius,
and the valuation of \(\varpi\). These classes exhaust
\(\mathbb Q/\mathbb Z\), proving compatibility of the invariants.
\end{proof}

\begin{lemma}[Deleting a nonempty finite set]
\label{lem:do27-open}
One has \(H^3(U,\mathbb G_m)=0\). For every
\(a\in H^2(U,\mathbb G_m)\),
\[
 \sum_{v\in S}\operatorname{inv}_v(a|_{K_v})=0.
\]
\end{lemma}

\begin{proof}
For \(j_U:U\hookrightarrow X\), the local direct-image input gives
\(R^{>0}j_{U*}\mathbb G_m=0\). The divisor sequence is
\[
 0\to\mathbb G_{m,X}\to j_{U*}\mathbb G_{m,U}
    \xrightarrow{(v)_{v\in S}}
       \bigoplus_{v\in S}i_{v*}\mathbb Z\to0.
\]
Since \(H^2(k_v,\mathbb Z)=\mathbb Q/\mathbb Z\) and
\(H^3(k_v,\mathbb Z)=0\), its long exact sequence contains
\[
 H^2(U,\mathbb G_m)\longrightarrow
 \bigoplus_{v\in S}\mathbb Q/\mathbb Z
 \xrightarrow{\delta_S}H^3(X,\mathbb G_m)
 \longrightarrow H^3(U,\mathbb G_m)\longrightarrow0.
\]
The first arrow is the local invariant map by the local input and
valuation computation. Naturality from the full divisor sequence gives
\(\operatorname{tr}_F\delta_S=\sum_{v\in S}\).
This sum is surjective because \(S\neq\varnothing\). The assertions
now follow from exactness.
\end{proof}

\subsection{A cochain model on the small site}

For an abelian sheaf \(\mathcal F\), we specify the cochains
to fix the maps and the cup products.
Choose a conservative family of geometric points of \(X\).
Let \(p^*\) be the family of stalk functors and \(p_*\) its right
adjoint. Put \(T=p_*p^*\). The unit and counit define the augmented
cosimplicial resolution
\[
 \mathcal F\longrightarrow T\mathcal F
       \rightrightarrows T^2\mathcal F\rightrightarrows T^3\mathcal F
       \longrightarrow\cdots.
\]
Use its unnormalized cochain complex \(I^q\mathcal F=T^{q+1}\mathcal F\),
with differential \(\sum_{i=0}^{q+1}(-1)^id^i\).
For an \'etale \(V\to X\), put
\(C^q(V,\mathcal F)=\Gamma(V,I^q\mathcal F)\).
Every geometric stalk of the augmented resolution has the extra
codegeneracy induced by the counit. Thus it is exact.
The terms are acyclic for sections over any such \(V\).
Indeed, the sections of \(p_*A\) are products of the groups \(A\)
indexed by lifts of geometric points to \(V\). This functor is exact.
Resolve \(A\) by injectives in the product category of abelian groups.
Since \(p^*\) is exact, \(p_*\) preserves injectives, proving the
acyclicity. This is the Godement construction
\cite{do27-Godement}; the argument here applies to the specified site.

Choose affine separated \'etale neighborhoods \(V\to X\) of \(v\)
with fiber \(\operatorname{Spec}k_v\) at \(v\), and with no fiber
at the other points of \(S\). Define
\[
 C(D_v,\mathcal F)=\varinjlim_V C(V,\mathcal F),\qquad
 C(K_v,\mathcal F)=\varinjlim_V C(V\cap U,\mathcal F).
\]
The limits of these schemes are \(D_v\) and \(\operatorname{Spec}K_v\).
The continuity theorem for \'etale cohomology identifies the cohomology
of these complexes with the indicated small-site groups
\cite[Theorem 59.51.3]{do27-Limits}.
Its hypotheses hold: the schemes are quasi-compact and separated,
and the transition morphisms can be taken affine. Filtered colimits
of abelian groups are exact. The same complexes apply to the restriction
of any sheaf on an open neighborhood containing \(v\).

Write \(C_U=C(U,\mathcal F)\),
\(C_D=\prod_{v\in S}C(D_v,\mathcal F)\), and
\(C_K=\prod_{v\in S}C(K_v,\mathcal F)\).
Restrictions define \(r:C_U\to C_K\) and \(s:C_D\to C_K\).
The finite gluing complex is
\begin{equation}
\label{eq:do27-complex}
 \begin{split}
 A_S^q(\mathcal F)&=C_U^q\oplus C_D^q\oplus C_K^{q-1},\\
 D(u,a,b)&=(du,da,ru-sa-db).
 \end{split}
\end{equation}
Complexes are zero in negative degrees. Equivalently,
\(A_S=\operatorname{Cone}(r-s)[-1]\) with the displayed sign
convention. The map from global cochains is
\[
 J_S:C(X,\mathcal F)\longrightarrow A_S(\mathcal F),
 \qquad z\longmapsto(z|_U,(z|_{D_v})_v,0).
\]

\begin{proposition}[Finite \'etale gluing]
\label{prop:do27-gluing}
The map \(J_S\) is a quasi-isomorphism. For every bilinear sheaf
pairing \(\mathcal F\otimes\mathcal G\to\mathcal H\), it respects
the cohomological cup pairing.
\end{proposition}

\begin{proof}
For every term \(I^q\mathcal F\), restriction from \(V\) to an open
subscheme of \(V\) is surjective. It is a projection between the
products indexed by lifts of geometric points. Consequently
\[
 0\longrightarrow C(X,\mathcal F)
 \longrightarrow C_U\oplus C_D
 \xrightarrow{r-s}C_K\longrightarrow0
\]
is degreewise exact. To check its kernel, represent the finitely many
local germs on neighborhoods \(V_v\). Shrink them until equality
with the given section on \(U\) holds on every \(V_v\cap U\).
The family \(U\amalg\coprod_vV_v\to X\) is an \'etale cover.
An overlap for distinct \(v\)'s lies over \(U\). On a self-overlap,
the complement of the diagonal also lies over \(U\): the chosen
closed fiber is a single residue-field isomorphism. Thus the sections
agree on every overlap and descend uniquely to \(X\).
Surjectivity on the right follows by lifting separately from
\(V_v\cap U\) to \(V_v\). Finiteness of \(S\) permits these
choices and refinements simultaneously.

For completeness, let \((u,a,b)\) be a cocycle. Lift \(b\) to
\(\ell\in(C_U\oplus C_D)^{q-1}\) with \((r-s)\ell=b\).
Subtracting \(D(\ell,0)\) leaves \(((u,a)-d\ell,0)\), a
cocycle in the kernel. This proves surjectivity of \(H(J_S)\).
If a kernel cocycle is \(D(\ell,b)\), lift \(b\) one degree lower
to \(m\). Then \(\ell-dm\) lies in the kernel and has the same
differential. This proves injectivity.

The functor \(T\) has a natural pairing: pair sections point by
point. It induces pairings on all cosimplicial degrees. For cochains
\(x\) and \(y\) of degrees \(p\) and \(q\), use
\[
 x\smile y=
 (d^{p+q}\cdots d^{p+1}x)\,\cdot\,((d^0)^py).
\]
The dot denotes the given coefficient pairing. This
Alexander--Whitney formula has
\(d(x\smile y)=dx\smile y+(-1)^px\smile dy\).
The cosimplicial identities cancel the interior cofaces in pairs;
the two remaining groups are exactly the two displayed summands.
It represents the derived cup pairing because the augmentation
restricts it to the coefficient pairing in degree zero.
Restrictions commute with every coface and with the pairing.

For \(x=(u,a,b)\in A_S^p(\mathcal F)\) and
\(y=(u',a',b')\in A_S^q(\mathcal G)\), set
\begin{equation}
\label{eq:do27-cup}
 x\smile y=
 \bigl(u\smile u',\ a\smile a',\
        b\smile ru'+(-1)^p sa\smile b'\bigr).
\end{equation}
Its last component has degree \(p+q-1\). The identity
\[
 ru\smile ru'-sa\smile sa'
 =(ru-sa)\smile ru'+sa\smile(ru'-sa')
\]
and the preceding Leibniz rule give
\(D(x\smile y)=Dx\smile y+(-1)^px\smile Dy\).
For associative coefficient pairings, the three terms in a triple
product are
\[
 b\smile ru'\smile ru'',\quad
 (-1)^p sa\smile b'\smile ru'',\quad
 (-1)^{p+q}sa\smile sa'\smile b''.
\]
They are the same under both associations. Finally \(J_S\) has last
component zero, so it commutes with this cup formula on cochains.
\end{proof}

For \(M\), multiplication in \(\mathbb Z/n\) gives an associative
cup product. For \(\mathbb G_m\) we use only specified bilinear
pairings, such as evaluation
\(M\otimes\underline{\operatorname{Hom}}(M,\mathbb G_m)\to\mathbb G_m\).
Multiplication of units is not a bilinear map
\(\mathbb G_m\otimes\mathbb G_m\to\mathbb G_m\).

\subsection{The degree-three trace and its sign}

\begin{theorem}[Trace by local primitives]
\label{thm:do27-trace}
Let \((u,a,b)\in Z^3 A_S(\mathbb G_m)\). Choose
\[
 h\in C^2(U,\mathbb G_m),\quad
 k_v\in C^2(D_v,\mathbb G_m),\qquad dh=u,\quad dk_v=a_v.
\]
Then
\begin{equation}
\label{eq:do27-trace}
 T_S[(u,a,b)]
   =\sum_{v\in S}\operatorname{inv}_v
           [\,r_vh-s_vk_v-b_v\,]
\end{equation}
is independent of the choices and the representative. It defines a
homomorphism \(H^3 A_S(\mathbb G_m)\to\mathbb Q/\mathbb Z\), and
\[
 T_S\circ H^3(J_S)=\operatorname{tr}_F.
\]
\end{theorem}

\begin{proof}
The primitives exist by Lemmas~\ref{lem:do27-local} and
\ref{lem:do27-open}. The cone cocycle equation gives
\(db=ru-sa\), so every cochain inside brackets is closed.
Changing \(k_v\) changes that bracket by the restriction of a
closed degree-two cochain on \(D_v\). Its class is zero by
Lemma~\ref{lem:do27-local}. Changing \(h\) adds restrictions of a
degree-two global cocycle on \(U\). Its local invariants sum to zero
by Lemma~\ref{lem:do27-open}.

If \((u,a,b)\) changes by \(D(p,q,t)\), use \(h+p\) and \(k+q\).
The bracket changes by \(dt\). Thus the value depends only on the
cohomology class. Add primitives to check additivity.

We now compare connecting maps, including their signs. Put
\(B=C_U\oplus C_D\), \(L=C_K\), and
\(C_X=C(X,\mathbb G_m)\). In the degreewise exact
sequence \(0\to C_X\to B\xrightarrow{r-s}L\to0\), lift a
closed \(t\in L^2\) to \(\ell\in B^2\).
The usual connecting class is represented by \(d\ell\in C_X^3\).
But
\[
 D(\ell,0)=(d\ell,t),\qquad
 [(0,0,t)]=-H^3(J_S)[\delta_{\mathrm{glue}}t].
\]
Hence \((0,0,-t)\) represents the positive connecting class.

Under valuation, \(H^2(K_v,\mathbb G_m)\) identifies with
\(H^2(k_v,\mathbb Z)\). We verify that the gluing boundary gives
the divisor boundary with its positive sign. Resolve the finite divisor
sequence by the same exact Godement functors. Taking sections gives
a degreewise exact sequence of complexes: every term of the
resolution is a product of stalk groups. These complexes also
compute the open and local generic groups, by the higher-direct-image
vanishing in the arithmetic input. Common double resolutions compare
them with the complexes already chosen.

Represent a residue class by a cocycle \(t\) of degree two, and
lift it globally to a meromorphic cochain \(\ell\).
Then \(d\ell=z\) is a \(\mathbb G_m\)-cocycle of degree three,
representing \(\delta_S[t]\). Its restriction to \(U\) supplies
\(h\), with \(dh=z_U\). At \(v\), choose a closed meromorphic
cochain \(\beta_v\) with valuation exactly \(t_v\).
Such a cochain exists: valuation is an isomorphism on degree-two
cohomology, and a discrepancy which is a coboundary can be removed
by lifting its primitive through the degreewise surjective valuation
map. The difference \(\ell|_{D_v}-\beta_v\) has valuation zero,
so it is a unit cochain \(k_v\). It satisfies \(dk_v=z_{D_v}\).
In the local generic complex,
\[
 r_vh-s_vk_v=\beta_v.
\]
Thus \((h,k)\) is a lift of \(\beta\) through \(r-s\), and
its differential is \(J_Sz\). This proves that
\(\delta_{\mathrm{glue}}[\beta]=\delta_S[t]\).
The reasoning takes place in degreewise exact resolutions and so
also proves the statement after their comparison maps. By~\eqref{eq:do27-input},
\[
 \operatorname{tr}_F\delta_{\mathrm{glue}}(\beta)
     =\sum_v\operatorname{inv}_v(\beta_v).
\]

Subtract \(D(h,k,0)\) from the original cocycle. The result is
\((0,0,b-rh+sk)\). The connecting-map computation now gives
exactly~\eqref{eq:do27-trace} and the asserted comparison.
\end{proof}

The sign is already visible over \(\mathbb F_3\). Let the complex
\(B\) have \(B^2=B^3=\mathbb F_3\) and differential the identity.
Let \(L\) be concentrated in degree two, with \(L^2=\mathbb F_3\),
and let \(r:B\to L\) be the identity in degree two. Its kernel
\(K\) is concentrated in degree three. Then
\[
 H^3(\operatorname{Cone}(r)[-1])
   =\mathbb F_3^2/\{(t,t):t\in\mathbb F_3\}.
\]
The positive connecting class of \(1\in L^2\) is represented by
\((1,0)\), so the normalized functional is \([(x,y)]\mapsto x-y\).
It sends the bare cone class \((0,1)\) to \(-1=2\pmod3\).
Reversing this sign sends that same class to \(1\), which contradicts
the positive connecting normalization.

For a single prime and a local class of invariant \(1/3\), the
cone class \((0,0,\alpha)\) has trace \(-1/3\).
Taking \(n=3\), \(F=\mathbb Q(\zeta_3)\), and the prime
above \(3\) realizes this case by~\eqref{eq:do27-carry}. Kummer theory gives a
class \(\gamma\in H^2(K_v,M)\) mapping to \(\alpha\), so the
finite-coefficient cone class also has trace \(-1/3\).
The opposite sign would disagree with the divisor orientation.
This case also retains a coefficient prime in the integer scheme.

\subsection{Passage to completed local fields}

We next replace the henselian fields by the local fields in the
finite adelic product \(\prod_{v\in S}F_v\).
This is a change of local cochain complexes, rather than a completion
of the coefficient group.

\begin{lemma}[Finite coefficients and completion]
\label{lem:do27-completion}
For \(K=K_v\), base change to \(\widehat K=F_v\) gives an
equivalence of finite \'etale field algebras. Thus the induced map of
absolute Galois groups is an isomorphism, after choosing compatible
separable closures. It preserves inertia and arithmetic Frobenius
on the unramified quotient.
\end{lemma}

\begin{proof}
We give the approximation argument. A henselian valuation extends
uniquely to an algebraic closure, so its automorphisms are isometries.
For a separable element \(\alpha\), put
\(c=\min_{\alpha'\neq\alpha}|\alpha-\alpha'|>0\), where
\(\alpha'\) runs through its conjugates. If \(|\alpha-\beta|<c\),
every automorphism fixing \(\beta\) fixes \(\alpha\).
Indeed, an automorphism \(\sigma\) fixing \(\beta\) satisfies
\[
 |\sigma\alpha-\alpha|
 \leq\max\{|\sigma\alpha-\beta|,|\beta-\alpha|\}
 =|\alpha-\beta|<c.
\]
It follows that \(K(\alpha)\subset K(\beta)\) when both elements
are algebraic over the henselian field under consideration.
If \(\alpha\in\widehat K\) is separable algebraic over \(K\),
choose \(\beta\in K\) sufficiently close to it. Apply the same
argument in the algebraic extension containing \(\alpha\).
It gives \(\alpha\in K\). Thus \(K\) is separably closed in
\(\widehat K\).

For a finite Galois extension \(E/K\), this implies linear
disjointness from \(\widehat K\). In fact
\(E\cap\widehat K=K\), and restriction identifies
\(\operatorname{Gal}(E\widehat K/\widehat K)\) with the subgroup
of \(\operatorname{Gal}(E/K)\) whose fixed field is that
intersection. The subgroup is the full group. Applying this to
Galois closures preserves degrees and all embeddings of finite
separable extensions.

Conversely, take \(\widehat E=\widehat K(\alpha)\), a finite
separable extension, and its monic minimal polynomial \(f\).
Scale \(\alpha\) to be integral. Approximate the coefficients of
\(f\) by coefficients in \(K\), obtaining a monic polynomial
\(g\) of the same degree. They can be chosen so that
\(|g(\alpha)|<|g'(\alpha)|^2\), and so that the resulting
Hensel root \(\beta\in\widehat E\) is closer to \(\alpha\)
than all its other conjugates. The estimate
\(|\beta-\alpha|\leq|g(\alpha)/g'(\alpha)|\) makes this explicit.
The preceding isometry argument over \(\widehat K\) gives
\(\widehat K(\alpha)\subset\widehat K(\beta)\).
Both fields lie in \(\widehat E\), so equality holds.
The degree of \(K(\beta)\) is at most \(\deg g\) and at least
\([\widehat E:\widehat K]\). These are the same number.
Thus \(K(\beta)\otimes_K\widehat K=\widehat E\).
Products and the embedding statement prove the equivalence.
Approximating coordinates in a finite field basis shows that each
finite extension is dense in its completed extension. A nonzero element
can therefore be approximated with strictly smaller error than its
absolute value. The approximant has the same valuation and, for units,
the same residue. Thus value groups and residue fields are unchanged.
The unramified subcategories and their Frobenius automorphisms agree.
\end{proof}

For clarity, the cochain map under this equivalence is not implicit.
The small \'etale topos of a field \(L\) is the topos of discrete
continuous \(G_L\)-sets. Finite \'etale objects give its finite
continuous sets, and sheaf descent extends this description to all
objects \cite[Lemmas 59.59.1--59.59.2]{do27-Field}.
Its standard cochains with a discrete module \(N\) are
\(C^q_{\mathrm{cont}}(G_L,N)\), with differential
\[
\begin{split}
 (df)(g_1,\ldots,g_{q+1})={}&g_1f(g_2,\ldots,g_{q+1})\\
 &+\sum_{i=1}^q(-1)^if(g_1,\ldots,g_ig_{i+1},\ldots,g_{q+1})\\
 &+(-1)^{q+1}f(g_1,\ldots,g_q).
\end{split}
\]
For a group map \(\theta:G_{F_v}\to G_{K_v}\) and a compatible
module map \(\lambda:N\to N'\), completion sends
\[
 f\longmapsto\bigl[(g_1,\ldots,g_q)\mapsto
       \lambda f(\theta g_1,\ldots,\theta g_q)\bigr].
\]
For \(M\), \(\lambda\) is the identity. For \(\mathbb G_m\),
it is the chosen inclusion of separable closures on their units.
For the latter coefficients, the map in degree zero is not an
isomorphism. The local invariant comparison concerns degree two.
The group cup formula is
\[
 (f\smile g)(g_1,\ldots,g_{p+q})=
 f(g_1,\ldots,g_p)\cdot
 (g_1\cdots g_p)g(g_{p+1},\ldots,g_{p+q}),
\]
for a specified equivariant bilinear pairing. Substitution shows
that the completion map respects this formula and the differential.
The bar resolution and the Godement resolution are acyclic resolutions
of the same coefficient object. Their common double resolution
provides comparison quasi-isomorphisms compatible with restriction
and the derived cup pairing.

Here is also a direct description before passing to the derived
category. For a map of the local schemes, choose geometric points
with compatible images. The adjunction maps give
\(f^{-1}T\mathcal F\to T'f^{-1}\mathcal F\), by sending a section
to each of its target stalks. Iterate this map in every cosimplicial
degree and apply the coefficient map. Taking sections gives the
cochain maps. The unit and counit identities make them commute with
cofaces, pairings, and the restriction square. Apply these maps to
the filtered neighborhood complexes above.

Use these comparison maps to define
\[
 \widehat A_S(\mathcal F)=
 \operatorname{Cone}\left(
 C(U,\mathcal F)\oplus
 \prod_{v\in S}C(\widehat R_v,\mathcal F)
 \xrightarrow{r-s}\prod_{v\in S}C(F_v,\mathcal F)
 \right)[-1].
\]
There is a cochain map \(\kappa_S:A_S\to\widehat A_S\).
For \(M\), it is a quasi-isomorphism: the field terms agree by
Lemma~\ref{lem:do27-completion}, and the ring terms have the same
residue-field cohomology. The cone long exact sequence proves the
claim. For \(\mathbb G_m\), it induces an isomorphism in degree
three. Indeed, both degree-three groups are the cokernel of
\[
 H^2(U,\mathbb G_m)\longrightarrow
          \prod_{v\in S}H^2(K_v,\mathbb G_m)
 \quad\text{or}\quad
 H^2(U,\mathbb G_m)\longrightarrow
          \prod_{v\in S}H^2(F_v,\mathbb G_m).
\]
Lemmas~\ref{lem:do27-local} and \ref{lem:do27-open} justify this
cokernel description. Local invariant compatibility identifies the
arrows. Formula~\eqref{eq:do27-trace} defines \(\widehat T_S\)
on the completed complex and gives
\[
 \widehat T_S H^3(\kappa_S)=T_S.
\]

\subsection{Orientation, boundary choices, and finite torsors}

Write the already constructed arithmetic orientation as
\[
 \operatorname{or}_{F,n,\zeta}
 :=\operatorname{inv}_{F,n,\zeta}
 =\operatorname{tr}_F\circ H^3(\epsilon):
 H^3(X,M)\xrightarrow{\sim}\tfrac1n\mathbb Z/\mathbb Z.
\]
Proposition~\ref{prop:do26-orientation} proves this isomorphism on
\(X\), including its coefficient primes.
For \(n=1\), all constructions with \(M\) have zero cohomology.
The comparison just proved is the commutative diagram
\begin{equation}
\label{eq:do27-comparison}
\begin{CD}
 H^3(X,M) @>{H^3(\epsilon)}>> H^3(X,\mathbb G_m)\\
 @V{H^3(\kappa_SJ_S)}V{\sim}V @VV{H^3(\kappa_SJ_S)}V\\
 H^3\widehat A_S(M) @>{H^3(\epsilon)}>>
 H^3\widehat A_S(\mathbb G_m) @>{\widehat T_S}>> \mathbb Q/\mathbb Z.
\end{CD}
\end{equation}
Composing either path with the trace gives
\(\operatorname{or}_{F,n,\zeta}\).
For a sheaf homomorphism \(f:M\to\mathbb G_m\), the same diagram
is the cup evaluation of degree three against degree zero.
Proposition~\ref{prop:do27-gluing} proves this compatibility for all
specified cup pairings, not only for \(f=\epsilon\).

To state the ramified boundary precisely, put
\[
 C_{c,S}(U,\mathcal F)=
 \operatorname{Cone}\left(C(U,\mathcal F)
         \xrightarrow r\prod_{v\in S}C(F_v,\mathcal F)\right)[-1],
 \quad D(u,b)=(du,ru-db).
\]
The subscript denotes this displayed finite cone. For \(\mathcal F=M\)
or \(\mathbb G_m\), its degree-three cohomology has a canonical
map to \(H^3(X,\mathcal F)\), obtained from
\((u,b)\mapsto(u,0,b)\) and the gluing comparison.
For \(\mathcal F=M\), this map is an isomorphism: the quotient
complex consists of the local ring terms, whose degrees two and
three vanish. For \(\mathbb G_m\), the same conclusion follows
from their positive-degree vanishing and the degree-three comparison.
On \(M\)-valued cochains the compact trace is
\begin{equation}
\label{eq:do27-compact-trace}
 \operatorname{Tr}_{c,S}[(u,b)]
  =\sum_{v\in S}\operatorname{inv}_v[\,r_vh-\epsilon b_v\,],
 \qquad dh=\epsilon u.
\end{equation}
In particular the value is in \(\tfrac1n\mathbb Z/\mathbb Z\).
This follows by its comparison with \(H^3(X,M)\), rather than an
assertion that each auxiliary \(\mathbb G_m\)-cochain is killed by \(n\).

\subsection{The positive boundary torsor and its dual}
\label{sec:do28-dual}

Put \(\Lambda=\tfrac1n\mathbb Z/\mathbb Z\) and
\[
 B_S=\prod_{v\in S}H^2(F_v,M),\qquad
 \sigma_\zeta:B_S\longrightarrow\Lambda,
 \quad(\eta_v)\longmapsto
       \sum_v\operatorname{inv}_v(\epsilon_*\eta_v).
\]
For \(u\in Z^3C(U,M)\), define
\[
 \mathcal H_S(u)=
 \{b\in\prod_{v\in S}C^2(F_v,M):db=ru\}/
                \prod_{v\in S}dC^1(F_v,M).
\]
It is a nonempty \(B_S\)-torsor by Lemma~\ref{lem:do27-local}.
For two points \(b,b'\), the notation \(b'-b\) means their unique
difference in \(B_S\). All torsor actions below are additive right
actions. The finite torsors carry the discrete topology.

The positive and negative pushouts are the \(\Lambda\)-torsors
\begin{equation}
\label{eq:do28-pushouts}
 L_S^\pm(u)=(\mathcal H_S(u)\times\Lambda)/\sim_\pm,
 \qquad (b+\eta,t)\sim_\pm(b,t\pm\sigma_\zeta\eta).
\end{equation}
Their actions are \([b,t]_\pm+a=[b,t+a]_\pm\).
For any fixed \(b\), the map \(t\mapsto[b,t]_\pm\) is a bijection.
Indeed, freeness of \(\mathcal H_S(u)\) gives a unique difference
between any two representatives. Thus the displayed actions are free
and transitive. The natural map \(b\mapsto[b,0]_+\) has positive
translation \(\sigma_\zeta\).

After the field-cochain identification, \(\mathcal H_S\) is Kim's
primitive torsor, and \(L_S^+\) is its positive sum pushout
\cite[\S2, p.~5, the definitions of \(H(\rho_S)\) and
\(L(\rho_S)\)]{Kim2016}. A scalar functional with negative translation
instead defines an equivariant map on \(L_S^-\).

For a \(\Lambda\)-torsor \(L\), define its opposite \(L^{\mathrm{op}}\)
to have the same underlying set and action \(\ell*a=\ell-a\).
Define its affine dual by
\[
 L^\vee=\{f:L\to\Lambda:f(\ell+a)=f(\ell)+a\},
 \qquad(f+a)(\ell)=f(\ell)+a.
\]
This is a \(\Lambda\)-torsor. Each \(\ell_0\in L\) determines the
map \(f_{\ell_0}(\ell)=\ell-\ell_0\). Every element of \(L^\vee\)
has this form, since every equivariant map has exactly one zero.

\begin{proposition}[Opposite torsor and dual evaluation]
\label{prop:do28-dual}
There are canonical equivariant isomorphisms
\[
 L_S^-(u)\xrightarrow{\iota_S}(L_S^+(u))^{\mathrm{op}}
       \xrightarrow{\sim}(L_S^+(u))^\vee.
\]
They are given by
\begin{equation}
\label{eq:do28-dual-map}
 \iota_S[b,t]_-=[b,-t]_+,
 \qquad
 \Psi_S([b,t]_-)([b',s]_+)
       =s+t+\sigma_\zeta(b'-b).
\end{equation}
In particular
\(E_S(\ell,m)=\Psi_S(m)(\ell)\) is a canonical evaluation pairing.
It satisfies
\(E_S(\ell+a,m+a')=E_S(\ell,m)+a+a'\).
\end{proposition}

\begin{proof}
The negative pushout relation maps under \(\iota_S\) to
\[
 [b+\eta,-t]_+=[b,-t+\sigma_\zeta\eta]_+.
\]
This is the image of \([b,t-\sigma_\zeta\eta]_-\), so the map
is well defined. Adding \(a\) on its source subtracts \(a\) in the
positive torsor, which is the opposite action.
The map \(L^{\mathrm{op}}\to L^\vee\), \(\ell_0\mapsto f_{\ell_0}\),
is equivariant because \(f_{\ell_0-a}=f_{\ell_0}+a\).
For \(\ell_0=[b,-t]_+\), its value at \([b',s]_+\) is the
right side of~\eqref{eq:do28-dual-map}.
This proves the formula and both isomorphisms. The evaluation identity
follows directly from addition in the two coordinates.
\end{proof}

This also gives a literal duality of complex lines. Let
\(\chi(a)=\exp(2\pi ia)\) and define
\(\mathcal L(L)=(L\times\mathbb C)/((\ell+a,z)\sim(\ell,\chi(a)z))\),
with its vector-space structure and usual complex topology defined in
any torsor coordinate.
The formula
\[
 ([\ell,z],[m,w])\longmapsto zw\chi(E_S(\ell,m))
\]
is independent of representatives and is a nonzero bilinear pairing.
It identifies \(\mathcal L(L_S^-)\) with
\(\mathcal L(L_S^+)^*\). These are finite-dimensional associated
lines, with no determinant or analytic trace asserted.

\begin{lemma}[A global primitive with finite coefficients]
\label{lem:do28-global-primitive}
Under the stated hypotheses, \(H^3(U,M)=0\).
For every \(a\in H^2(U,M)\), one has
\(\sigma_\zeta(r a)=0\).
\end{lemma}

\begin{proof}
The sequence in Lemma~\ref{lem:do27-open}, together with
\(H^2(X,\mathbb G_m)=0\), identifies
\[
 H^2(U,\mathbb G_m)
 \xrightarrow{\sim}
 \ker\left(\bigoplus_{v\in S}\mathbb Q/\mathbb Z
                         \xrightarrow{\sum_v}\mathbb Q/\mathbb Z\right).
\]
Choose \(v_0\in S\). Projection to the other coordinates identifies
this kernel with \((\mathbb Q/\mathbb Z)^{|S|-1}\), the omitted
coordinate being the negative sum. Hence multiplication by \(n\)
on \(H^2(U,\mathbb G_m)\) is surjective.
Every prime above \(n\) lies in \(S\). The \'etale Kummer sequence
is therefore exact on \(U\), and \(e_\zeta:M\to\mu_n\) is an
isomorphism there. Its long exact sequence, with
\(H^3(U,\mathbb G_m)=0\), gives \(H^3(U,M)=0\).
Finally apply \(\epsilon_*\) to \(a\), and use the zero-sum
statement of Lemma~\ref{lem:do27-open}.
\end{proof}

\begin{theorem}[The global boundary point and the arithmetic trace]
\label{thm:do28-state}
For \(u\in Z^3C(U,M)\), choose \(\beta\in C^2(U,M)\) with
\(d\beta=u\). There is a well-defined point
\begin{equation}
\label{eq:do28-state}
 \lambda_S(u)=[r\beta,0]_+\in L_S^+(u),
\end{equation}
independent of \(\beta\). For every \(b\in\mathcal H_S(u)\),
\begin{equation}
\label{eq:do28-state-trace}
 \tau_u(b):=\operatorname{Tr}_{c,S}[(u,b)]
   =\sigma_\zeta(r\beta-b)
   =\lambda_S(u)-[b,0]_+.
\end{equation}
Consequently the map
\[
 \Theta_u:L_S^-(u)\longrightarrow\Lambda,
 \qquad [b,t]_-\longmapsto\tau_u(b)+t,
\]
is an equivariant isomorphism and satisfies
\(\Theta_u(m)=E_S(\lambda_S(u),m)\).
Its zero is the dual point \(\lambda_S(u)^\vee=\Psi_S^{-1}(f_{\lambda_S(u)})\).
\end{theorem}

\begin{proof}
Lemma~\ref{lem:do28-global-primitive} gives \(\beta\).
Two choices differ by a global two-cocycle \(a\), whose local
invariants sum to zero. The positive pushout relation gives
\([r\beta+r a,0]_+=[r\beta,\sigma_\zeta(r a)]_+=[r\beta,0]_+\).
Thus \(\lambda_S(u)\) is independent of that choice.

Take \(h=\epsilon\beta\) in~\eqref{eq:do27-compact-trace}.
It gives the first equality in~\eqref{eq:do28-state-trace}.
The positive pushout relation gives
\([r\beta,0]_+=[b,\sigma_\zeta(r\beta-b)]_+\), which proves
the torsor difference formula. In particular
\(\tau_u(b+\eta)=\tau_u(b)-\sigma_\zeta\eta\).
This is exactly the relation needed to define \(\Theta_u\).
It is equivariant and therefore is an isomorphism of torsors.
Substitute \(\lambda_S(u)=[b,\tau_u(b)]_+\) in the evaluation
formula to obtain \(\Theta_u(m)=E_S(\lambda_S(u),m)\).
In coordinates its zero is \([b,-\tau_u(b)]_-\), independent of
\(b\). Under \(\Psi_S\) this is the function with zero
\(\lambda_S(u)\), as asserted.
\end{proof}

Kim's global boundary point is \([r\beta,0]_+\)
\cite[\S2, p.~6, the displayed definition of \(\operatorname{CS}_c\)]{Kim2016}.
Theorem~\ref{thm:do28-state} identifies the cone trace with its
difference from the chosen boundary point. Kim also prints the scalar
coordinate \(\sum_v\operatorname{inv}_v(b_v-r_v\beta)\)
\cite[\S5, p.~12]{Kim2016}. For the same cochains and local
invariants that displayed coordinate equals \(-\tau_u(b)\).
Negation is the comparison between these two scalar coordinates.
The divisor normalization remains~\eqref{eq:do27-input}.

\subsection{Gluing and orientation compatibility}
\label{sec:do28-gluing}

For two \(\Lambda\)-torsors define
\(L_1\otimes_\Lambda L_2=(L_1\times L_2)/((\ell_1+a,\ell_2)\sim
(\ell_1,\ell_2+a))\), with addition in either factor.
Write \(\sigma_T\) for the sum of local invariant maps over a
finite subset \(T\), with the root \(\zeta\) fixed.
A finite partition \(S=S_1\amalg S_2\) gives canonical isomorphisms
\begin{equation}
\label{eq:do28-block-gluing}
 L_S^\pm(u)\cong L_{S_1}^\pm(u)\otimes_\Lambda L_{S_2}^\pm(u),
 \qquad
 ([b_1,t_1]_\pm,[b_2,t_2]_\pm)\longmapsto
       [(b_1,b_2),t_1+t_2]_\pm,
\end{equation}
where the arrow displays the inverse to the indicated isomorphism.
The factors mean boundary torsors for restrictions of this same \(u\).
Indeed, \(B_S=B_{S_1}\times B_{S_2}\) and
\(\sigma_S=\sigma_{S_1}+\sigma_{S_2}\). Substitution proves that
the formula respects both pushout relations and the contracted
product relation. It is equivariant and hence bijective.
The evaluation formula gives
\[
 E_S(\ell_1\otimes\ell_2,m_1\otimes m_2)
   =E_{S_1}(\ell_1,m_1)+E_{S_2}(\ell_2,m_2).
\]
It follows that this gluing commutes with the opposite and dual
identifications. It is associative under further finite partitions,
because addition of the displayed coordinates is associative.
For an empty block, use the torsor \(\Lambda\) with point zero as
the tensor unit.

A chosen global primitive \(\beta\) gives block points
\(\lambda_{S_i}(\beta)=[r_{S_i}\beta,0]_+\).
Their tensor is \(\lambda_S(u)\).
Changing \(\beta\) by a global cocycle \(a\) translates those
block points by \(\sigma_{S_i}(r a)\). Their sum is zero by
Lemma~\ref{lem:do28-global-primitive}. Thus the glued point is
independent of \(\beta\), even when the individual block points
are not.

There is also a direct gluing formula for two degree-three cocycles.
Suppose \(u_1,u_2\in Z^3C(U,M)\) satisfy \(r u_1=r u_2\) as
cochains. Their boundary torsors have the same definition, so write
\(L^+\) for this common torsor. Then
\begin{equation}
\label{eq:do28-difference-gluing}
 E_S(\lambda_S(u_1),\lambda_S(u_2)^\vee)
   =\operatorname{Tr}_{c,S}[(u_1-u_2,0)].
\end{equation}
To prove the formula, take a common boundary primitive \(b\).
The left side is \(\lambda_S(u_1)-\lambda_S(u_2)\), which equals
\(\tau_{u_1}(b)-\tau_{u_2}(b)\) by
Theorem~\ref{thm:do28-state}. Subtraction of the two relative
cocycles gives \((u_1-u_2,0)\), proving the equality.
Through the finite \'etale gluing isomorphism, the right side is the
arithmetic orientation of the corresponding class in \(H^3(X,M)\).
Thus dual evaluation has exactly the difference sign required by
that orientation.

Let \(A\) be a finite group, and let \(P_X\) be a finite \'etale
right \(A\)-torsor on \(X\). Fix \(c\in H^3(A,\mathbb Z/n)\)
with trivial coefficient action. The exact associated-sheaf functor
\(N\mapsto P_X\times^A\underline N\) acts on Yoneda extensions
and gives \(f_{P_X}^*c\in H^3(X,M)\).
Choose a cocycle \(z\) representing this class and put \(u=z_U\).
At each deleted prime choose \(a_v\) with
\(da_v=z|_{\widehat R_v}\). Local ring cohomology makes the
class of \(a_v\) modulo coboundaries unique. Put \(b^0=sa\).
Subtracting \(D(0,a,0)\) from \((z_U,z_D,0)\) gives
\((z_U,0,b^0)\). The established trace comparison therefore gives
\begin{equation}
\label{eq:do28-unramified-orientation}
 \operatorname{or}_{F,n,\zeta}(f_{P_X}^*c)
   =\lambda_S(u)-[b^0,0]_+
   =E_S(\lambda_S(u),[b^0,0]_-).
\end{equation}
Notice that \([b^0,0]_-\) is the dual of the positive boundary
point \([b^0,0]_+\). It need not be the dual of the global point
\(\lambda_S(u)\). This distinction permits a nonzero orientation.
For a torsor on \(U\) that does not extend to \(X\), the global
point \(\lambda_S(u)\) still exists. The scalar requires a boundary
point, and its translation remains \(-\sigma_\zeta\).

\subsection{Functoriality of the dual comparison}
\label{sec:do28-functoriality}

A boundary cocycle \(v\in\prod_S Z^3C(F_v,M)\) has its primitive
torsor \(\mathcal H(v)\). A boundary two-cochain \(q\) satisfying
\(dq=v'-v\) gives a transport
\[
 F_q:\mathcal H(v)\longrightarrow\mathcal H(v'),
 \quad[b]\longmapsto[b+q],\qquad
 F_q^\pm[b,t]_\pm=[b+q,t]_\pm.
\]
Replacing \(q\) by a coboundary leaves these maps unchanged.
Composition satisfies \(F_{q'}F_q=F_{q+q'}\).
Differences of two primitive points are unchanged by translation.
Thus~\eqref{eq:do28-dual-map} gives the exact identities
\begin{equation}
\label{eq:do28-functor-dual}
 \iota F_q^-=(F_q^+)^{\mathrm{op}}\iota,
 \qquad
 \Psi F_q^-(m)=\Psi(m)\circ(F_q^+)^{-1}.
\end{equation}
The inverse in the second formula specifies the variance of the dual.
In particular
\(E_{v'}(F_q^+\ell,F_q^-m)=E_v(\ell,m)\).
These identities apply to the local cochain translations used in
Kim's torsor functor \cite[\S2, pp.~5--6]{Kim2016}.
They require no choice of a torsor origin.

For global cocycles \(u,u'\), let \(u'-u=dz\).
The compatible boundary transport is \(q=rz\).
One has
\[
 (u',b+rz)-(u,b)=D(z,0),\quad
 F_{rz}^+\lambda_S(u)=\lambda_S(u'),\quad
 \Theta_{u'}F_{rz}^-=\Theta_u.
\]
The first identity proves trace compatibility. For the second, use
\(\beta+z\) as a primitive for \(u'\).
If another cochain \(z'\) has the same differential, then
\(z'-z\) is a global two-cocycle. Reciprocity gives
\(\sigma_\zeta(r(z'-z))=0\), so the two induced maps on
\(L_S^+\) and \(L_S^-\) agree.

This gives a precise functor on cocycles: its objects are the global
three-cocycles, and its arrows are two-cochains \(z\) modulo
\(dC^1(U,M)\), with \(dz=u'-u\). Composition is addition.
The constructions \(L_S^\pm\), \(\Psi_S\), \(\lambda_S\), and
\(\Theta\) respect this composition by the displayed identities.
Every automorphism acts trivially on the pushed-out torsors, since it
is a global two-cocycle. Different choices of representatives therefore
give canonically compatible torsor maps.

The associated-sheaf functor of a finite torsor is exact and preserves
tensor pairings, as can be checked on a cover trivializing the torsor.
An isomorphism of torsors gives a natural isomorphism of these functors
which is the identity on the trivial coefficient endpoints.
It therefore gives equal pulled-back Yoneda classes. Choose a two-cochain
\(z\) whose differential is the difference of the chosen
three-cocycles. The preceding independence from \(z\) makes the
maps on \(L_S^\pm\) independent of this choice. For a composite
of torsor isomorphisms, the sum of the two cochains has differential
equal to the difference for the composite. Thus the maps compose correctly.
The global point and its dual evaluation are natural under these maps.
Together with the cup comparison in Proposition~\ref{prop:do27-gluing},
this proves functorial compatibility with the orientation pairing.

We specify how this statement compares actual group cochains with
sheaf cochains. A finite continuous group cochain factors through a
finite quotient. A finite collection of cochains and their differential
identities therefore factors through one finite quotient \(Q\) of
\(\pi_1(U)\). The corresponding finite \'etale torsor gives the
exact functor \(T_Q\). Apply it to the coinduced bar resolution of
\(M\), and resolve its terms by the Godement functors. The resulting
first-quadrant double complex is a resolution of \(M\).
Its total differential is \(d_h+(-1)^p d_v\), with \(p\) the
horizontal degree. An invariant bar cochain maps to its associated
global section and then to Godement degree zero. The augmentation
makes its other differential zero, so this is a cochain map.
Passing through larger finite quotients gives the continuous-cochain
comparison. The same construction on each local field is the
field-site equivalence, so it is a quasi-isomorphism there.
All constructions use restriction of the finite torsor, so they give
compatible global and local maps through the common double resolutions.
No global \(K(\pi,1)\) assertion is needed.

For a square presented with a cochain homotopy, write its comparison
maps as \(F_U,F_\partial\) and fix a degree-minus-one map \(H\)
with
\[
 r'F_U-F_\partial r=dH+Hd.
\]
The induced primitive map is explicitly
\[
 (u,b)\longmapsto(F_Uu,F_\partial b+Hu).
\]
Its boundary equation follows by applying \(d\).
If \(d\beta=u\), then
\[
 F_\partial r\beta+Hu-r'F_U\beta=-dH\beta.
\]
Thus it carries \([r\beta,0]_+\) to the positive point represented
by \(r'F_U\beta\). It also intertwines \(\Psi\) and \(\Theta\),
since the local comparisons preserve the invariant maps.
This proves that Kim's global point, when represented by a continuous
group primitive, maps to~\eqref{eq:do28-state}. The argument needs
a comparison map, not a global cohomology isomorphism.

Finally replace \(\zeta\) by \(\zeta^a\), where
\(a\in(\mathbb Z/n)^\times\). Then
\(\sigma_{\zeta^a}=a\sigma_\zeta\) and
\(\tau_{u,\zeta^a}=a\tau_{u,\zeta}\).
The maps between the pushed-out torsors are
\[
 [b,t]_{\pm,\zeta}\longmapsto[b,at]_{\pm,\zeta^a}.
\]
They are equivariant for the automorphism \([a]\) of \(\Lambda\).
They carry the global points to the global points and multiply the
evaluation pairing by \(a\). The opposite action commutes with
\([a]\). Hence the dual comparison has the same primitive-root
normalization as the arithmetic orientation.

For a concrete sign test, take \(n=3\) and a boundary class \(\eta\)
with \(\sigma_\zeta\eta=1/3\). The positive boundary point changes
by \(+1/3\), its opposite changes by \(-1/3\), and
\(\tau_u(b+\eta)=\tau_u(b)-1/3\).
In the ramified example \(F=\mathbb Q(\zeta_9)\), with \(S\)
containing the primes above \(3\) and \(7\), take the torsor
\(y^3=7\) on \(U\). Its lift \(y^9=7\) is a \(C_9\)-torsor.
The ninth roots of unity lie in the unramified degree-three extension
of \(\mathbb Q_7\): \(9\mid7^3-1\), and Hensel's lemma lifts
their simple residue roots. Hence each prime over \(7\) in \(F\)
has \(v(7)=1\). Both displayed polynomials are Eisenstein there,
so these torsors are ramified at those primes. Use representatives
\(0,\ldots,8\) for \(C_9\), and write \(\bar i\) for the
residue of \(i\) in \(\{0,1,2\}\). On \(C_3\), let
\[
 c(i,j,k)=i\left\lfloor\frac{j+k}{3}\right\rfloor\pmod3.
\]
Its pullback to \(C_9\) has the primitive
\(B(i,j)=\bar i\lfloor j/3\rfloor\pmod3\).
Indeed, expanding \(dB\) cancels the quotients of \(j\) and \(k\),
leaving \(\bar i\lfloor(\bar j+\bar k)/3\rfloor\).
Pull back \(B\) to a global primitive \(\beta\).
For \(b=r\beta\), the global point is \([b,0]_+\) and the trace
is zero. Replacing one local primitive by \(b+\eta\), with invariant
\(1/3\), leaves the global point fixed and changes its coordinate
to \(-1/3\). This tests the opposite-torsor comparison on a fixed
ramified arithmetic torsor.

\subsection{Restriction and refinement of the support}

Let \(S\subset T\) be nonempty finite sets as above, and put
\(U_T=X\setminus T\). There is a cochain map
\(R_{ST}:A_S(\mathcal F)\to A_T(\mathcal F)\).
Restrict the global component from \(U\) to \(U_T\).
Keep the old local components. At \(w\in T\setminus S\), set
\[
 a_w=u|_{D_w},\qquad b_w=0.
\]
The square of restriction maps commutes, so this formula commutes
with \(D\). Formula~\eqref{eq:do27-cup} shows that it preserves
the cup pairing. Moreover
\[
 R_{ST}J_S=J_T,\qquad R_{TV}R_{ST}=R_{SV}\quad(S\subset T\subset V).
\]
The same formulas give maps for completed local terms.

The trace is compatible on cohomology. To see this directly, let
\((u,a,b)\) be a degree-three \(\mathbb G_m\)-cocycle and choose
\(dh=u\). At a new prime choose the local primitive
\(k_w=h|_{D_w}\). Its contribution to~\eqref{eq:do27-trace} is
\(r_wh-s_wk_w=0\). The old contributions are unchanged. Hence
\[
 T_T H^3(R_{ST})=T_S,
 \qquad \widehat T_T H^3(\widehat R_{ST})=\widehat T_S.
\]

For a degree-three class in the finite compact cone, first map it
to the gluing complex and then restrict. A new local ring component
is \(u|_{\widehat R_w}\). Choose \(a_w\) with
\(da_w=u|_{\widehat R_w}\) and remove it as above.
The new boundary primitive is \(b_w=s_wa_w\).
Its ambiguity is zero in cohomology because
\(H^2(\widehat R_w,M)=0\). This gives a canonical refinement map
in degree three, and it preserves~\eqref{eq:do27-compact-trace}.
For a torsor unramified at the new prime, these are precisely the
primitives supplied by its local ring extension.

The same refinement acts on the positive and negative boundary torsors.
For \(w\in T\setminus S\), the local primitive class
\([a_w]\), with \(da_w=u|_{\widehat R_w}\), is unique modulo
coboundaries. Its generic restriction is a distinguished boundary
point. Define
\[
 \mathcal R_{ST}^\pm[b,t]_\pm
       =[(b,(s_wa_w)_{w\in T\setminus S}),t]_\pm.
\]
If \(a_w\) changes by a closed ring cochain, its class in degree
two is zero. Thus the formula is independent of the chosen primitives.
On the old acting group, the new sum map is \(\sigma_S\), so the
formula respects both pushout relations. It is an equivariant
isomorphism of \(\Lambda\)-torsors.

For a global primitive \(\beta\) on \(U\), choose
\(a_w=\beta|_{\widehat R_w}\). This proves
\[
 \mathcal R_{ST}^+\lambda_S(u)=\lambda_T(u|_{U_T}),\qquad
 \Theta_{u|_{U_T}}\mathcal R_{ST}^-=\Theta_u.
\]
Moreover
\[
 \Psi_T\mathcal R_{ST}^-(m)
    =\Psi_S(m)\circ(\mathcal R_{ST}^+)^{-1}.
\]
The last identity follows by subtracting two refined primitive tuples:
their new components have difference zero. Hence the evaluation pairing
is unchanged. Refinements compose, since adjoining these unique local
ring primitive classes in either order gives the same tuple.
For \(u'=u+dz\), choose the new ring primitive \(a_w+z|_{\widehat R_w}\).
This also proves commutation with the global cocycle transports.
Thus support refinement commutes with the orientation functor, its
positive boundary point, and its opposite and dual torsors.

Changing an \'etale neighborhood representative refines the filtered
system defining \(C(D_v)\) and \(C(K_v)\). Equality in those
colimits means equality after a common refinement. Every formula
uses only finitely many components and finitely many cofaces.
It therefore descends to these colimits and commutes with such
refinement. A change of the cochain resolution gives the same
cohomological maps through the common double resolution and the
fixed connecting-map convention.

Local unramified primitive classes, the global boundary point, and a
prepared state have different carriers. The local primitives belong
to quotients of degree-two cochains on the local rings. The global
point \(\lambda_S(u)\) belongs to the finite
\(\Lambda\)-torsor \(L_S^+(u)\) and determines the nonzero vector
\([\lambda_S(u),1]\) in its associated complex line. Interpreting this
vector as a prepared state requires a specified quantum state space and
a preparation map compatible with that theory's gluing operations.
The finite torsor construction supplies no such quantum comparison.

The finite adelic cochain trace therefore supplies the oriented scalar
for an unramified finite torsor and the difference coordinate for a
ramified torsor with a chosen boundary point. A comparison
with integration of Schwartz--Bruhat functions, a Mellin transform,
or a determinant requires a map from the arithmetic coefficient
complex to that analytic carrier. No such map is part of the
finite gluing theorem. A continuous arithmetic Chern--Simons theory
would additionally require its continuous gauge group, coefficient
complex, topology, and compatible gluing and state maps. The finite
sums above imply no vanishing statement for an infinite product of
cohomology groups.
